% TOP_PROBLEM: {"id": 30006687, "title": "P versus NP ∩ coNP", "category_id": 15, "category": {"id": 15, "name": "computer_science", "display_name": "Computer Science", "description": "Computational complexity, algorithms, and theoretical CS.", "slug": "computer-science", "order_index": 15, "created_at": "2026-07-31T15:26:25.671Z"}, "status": "open"}
% FIRST_POSTED: 2026-09-27
% !TeX program = lualatex
% ATTEMPT_STATUS: unresolved
\documentclass[11pt]{article}
\usepackage[margin=25mm]{geometry}
\usepackage{fontspec}
\setmainfont{Latin Modern Roman}
\usepackage{amsmath,amssymb,mathtools}
\usepackage{unicode-math}
\setmathfont{Latin Modern Math}
\usepackage{xurl,hyperref}
\hypersetup{colorlinks=true,urlcolor=blue}
\setcounter{secnumdepth}{0}
\setlength{\parindent}{0pt}
\setlength{\parskip}{5pt}
\setlength{\emergencystretch}{3em}
\begin{document}
\section*{\texorpdfstring{P versus NP ∩ coNP}{P versus NP ∩ coNP}}\label{30006687}
\subsection{Definitions and mathematical statement}
A language lies in $\mathsf{NP}\cap\mathsf{coNP}$ when both membership and nonmembership have polynomial-length certificates verifiable in deterministic polynomial time. Thus there are verifiers $V_1,V_0$ and polynomial length bounds satisfying
\[
 x\in L\iff\exists y\,V_1(x,y)=1,\qquad
 x\notin L\iff\exists z\,V_0(x,z)=1.
\]
Determine whether every such language belongs to $\mathsf P$, that is, $\mathsf P=\mathsf{NP}\cap\mathsf{coNP}$. Algorithms are uniform and unrelativized, and inputs are total languages rather than promise problems.
\par\smallskip\textit{Formulation and concept references:} \href{https://www.scottaaronson.com/democritus/lec6.html}{[S1]}.

\subsection{Short English statement}
When both answers have efficiently checkable certificates, must the answer itself be efficiently computable?
\subsection{Research attempt}
\paragraph{Scope, source check, and method.}
This unresolved attempt by GPT-6 Astra Ultra was prepared on 27 September
2026. It examines the gap between verifying either answer and finding the
answer, using an explicit total factor-threshold language and a careful
analysis of certificate search. The two verifiers in the definition belong
to a fixed language; their correctness and coverage hold on every input.
No advice, oracle, or unmentioned promise is introduced.

Aaronson's lecture [S1] gives the class-comparison context. The recent
permutation paper [S3] explicitly obtains separations relative to random
permutations and related randomness notions. These do not constitute the
unrelativized separation requested here. The only substantial external
algorithm used in the concrete certificate construction below is the
deterministic polynomial-time primality test of Agrawal--Kayal--Saxena [E1].
That theorem recognizes prime numbers; it does not output their factors.
All reductions and certificate bookkeeping used here are proved explicitly.

The easy containment $\mathsf P\subseteq\mathsf{NP}\cap\mathsf{coNP}$
uses the deterministic decider itself as a verifier for each answer, with
an empty witness. To prove strict containment one needs a total language
having both certificate types but no deterministic polynomial-time decider.
To prove equality one needs a uniform decider for every fixed pair of
correct two-sided certificate relations. Neither task is accomplished by
exhibiting an interesting pair of verifiers alone.

\paragraph{Attack A: a fully total decision version of factoring.}
Fix a standard, efficiently parsed binary encoding of pairs of nonnegative
integers. Define
\begin{equation}\label{a34:factorlanguage}
 F=\{\langle N,K\rangle:N\ge2,\ \exists d\in\mathbb Z,
                         \ 2\le d<N,\ d\le K,\ d\mid N\}.
\end{equation}
Malformed strings and pairs with $N<2$ are not in $F$. No promise of
compositeness, semiprimality, or a restricted threshold is imposed. The bit
length of a valid input is comparable to the sum of the bit lengths of
$N$ and $K$, and polynomial means polynomial in that length.

\paragraph{Lemma 34.1: both answers have short certificates.}
The language $F$ belongs to $\mathsf{NP}\cap\mathsf{coNP}$.

\emph{Yes-certificates.} A witness is the divisor $d$. Verify the bounds
$2\le d<N$, $d\le K$, and compute $N\bmod d$. Its length is at most
$O(\log(N+2))$, and ordinary integer arithmetic verifies it in polynomial
time. No claim that $d$ itself is prime is needed.

\emph{No-certificates.} Malformed inputs, $N<2$, and $K<2$ can be recognized
and rejected by the language decider immediately, so their nonmembership
certificates are a fixed tag. In the remaining cases provide a list of
primes, repeated with multiplicity,
\[
                       p_1,\ldots,p_t,\qquad N=\prod_{i=1}^t p_i.
\]
Check the product exactly and check primality of each entry using [E1].
The no-verifier then accepts if either $t=1$ or every $p_i>K$.
The product condition with prime entries implies $t\le\log_2 N$, and the
list has polynomial encoded length, for example $O((\log(N+2))^2)$ with
simple delimiters.

If $t=1$, $N$ is prime and has no proper divisor. If $t\ge2$, every prime
factor is proper. A divisor $d\ge2$ of $N$ contains a prime factor $p_i$
with $p_i\le d$. Thus all prime factors being greater than $K$ precludes
an allowed divisor. Conversely, if $N$ is composite and has no allowed
divisor, none of its prime factors can be at most $K$; its complete prime
factorization is therefore an accepted no-certificate. This proves soundness
and completeness for both sides, including $K\ge N$ and prime $N$.
\hfill$\square$

The distinction between existence and discovery is visible here. A complete
factorization can be checked quickly even though no deterministic
polynomial-time construction of such a list is supplied. The no-verifier
must verify that the supplied prime factors multiply to all of $N$;
an incomplete list would not rule out a smaller omitted factor.

\paragraph{Lemma 34.2: deciding this language is equivalent to factoring.}
The following are equivalent: $F\in\mathsf P$, and there is a deterministic
polynomial-time algorithm outputting a complete prime factorization of every
binary integer $N\ge2$.

\emph{Proof, decision to factorization.} Given $N$, query $F$ at $(N,N-1)$.
A no answer says $N$ is prime. A yes answer says it is composite, and the
predicate $K\mapsto\mathbf1_F(N,K)$ is monotone. Binary search on
$2\le K\le N-1$ finds its smallest true threshold using $O(\log N)$
queries. That threshold is the smallest prime factor: the smallest proper
divisor cannot be composite, since a prime divisor of it would be a smaller
proper divisor of $N$. Output this prime, replace $N$ by its quotient, and
repeat until the quotient is one or is declared prime.

Every removal divides the current integer by at least two. There are at
most $\log_2 N$ removals, and every query has $O(\log N)$ bits. Thus
$O((\log N)^2)$ calls to a polynomial-time decider, together with polynomial
integer arithmetic, suffice. The search uses binary thresholds and never
loops through $N$ candidate integers.

\emph{Proof, factorization to decision.} First reject malformed or out-of-range
inputs. Factor $N$. If the list has one prime entry, answer no. Otherwise
compare the smallest prime factor with $K$. The equivalence established
in Lemma 34.1 proves this decides \eqref{a34:factorlanguage}. Polynomial
time is measured in the binary input length in both directions.
\hfill$\square$

Consequently equality of the two classes in the target would imply
deterministic polynomial-time factoring. A proof that this particular total
language is outside P would give a negative answer to the target. Neither
a lack of a known factoring algorithm nor success of quantum factoring is
such a classical lower bound. Conversely, a classical factoring algorithm
would settle this candidate language but would not establish the class
equality for all two-sided certificate relations.

\paragraph{Attack B: convert two certificates into a total search task.}
Pad the witness bounds to a common polynomial $p(n)$, using an explicit
length field if necessary. Define the polynomially checkable relation
\[
 R(x,b,w)\quad\Longleftrightarrow\quad
      b\in\{0,1\},\quad |w|\le p(|x|),\quad V_b(x,w)=1.
\]
For every input $x$, some pair $(b,w)$ satisfies $R$, and all satisfying
pairs have the same bit $b=1_L(x)$. A polynomial-time algorithm finding any
such pair would decide $L$. This gives a sufficient search route, but
totality alone does not supply the search algorithm.

Exhaustive search through both types of witness does terminate, with an
upper bound exponential in the polynomial witness length. If either side
has witnesses of length at most $c\log(n+2)$, exhaustive search on that
side alone is polynomial and decides the language: acceptance of one
witness gives that side's answer, and absence of all witnesses gives the
other. For a general polynomial bound $p(n)$, the same enumeration has
$2^{\Theta(p(n))}$ candidates and proves no polynomial runtime.

\paragraph{The prefix-extension gap in search-to-decision.}
A natural strategy is to ask whether a desired witness can begin with a
given prefix. For yes-certificates this predicate is
\[
 S=\{(x,u):\exists w\text{ extending }u,\ |w|\le p(|x|),\ V_1(x,w)=1\}.
\]
It is in NP. However a no answer need not be certified by $V_0$: there
may be a valid yes-certificate for $x$ beginning with another prefix.
For an explicit example take the trivial language of all strings and let
$V_1(x,w)$ accept exactly when $w=0^{|x|}$, while $V_0$ always rejects.
For $|x|\ge1$ the prefix $1$ has no extension, although no $V_0$ certificate
exists. Thus the proposed complement verifier fails on a concrete input.
This does not prove that $S$ is outside coNP; it proves that the original
pair of verifiers does not automatically certify prefix nonextendibility.

Accordingly, even assuming the equality asked in this problem would not
justify invoking it on every NP prefix predicate. Under the stronger
assumption $\mathsf P=\mathsf{NP}$, those queries could be answered and
ordinary prefix search would work. Importing that stronger assumption here
would change the problem.

\paragraph{Why finding the presented witness can be a stronger demand.}
Let $T(x,w)$ be any polynomial-time decidable, polynomially balanced relation
with at least one witness for every input. These are the total NP search
 relations commonly grouped under TFNP [E2]. Set $L=\{0,1\}^*$, let
$V_1(x,w)=T(x,w)$, and let $V_0$ always reject. This is a valid two-sided
certificate description of a language already in P. Yet finding a witness
for this particular $V_1$ is exactly the original total search problem.
Therefore a method that efficiently finds witnesses for \emph{every supplied}
valid verifier pair would solve all such total search problems, even for
constant-answer languages. The class-equality question only asks to decide
the answer; it does not demand efficient search for arbitrary presented
certificate systems. This construction identifies the additional obligation
in the proposed search route, without asserting a separation for TFNP.

There is a related promise trap in making the verifier pair itself part of
the input. Arbitrary pairs of polynomial-time programs need not be sound,
disjoint, or collectively total. Two always-rejecting verifiers leave every
input uncovered; two always-accepting verifiers certify both answers. A
universal problem restricted to well-behaved pairs therefore has a semantic
promise unless one supplies a separate, sound way to enforce that condition.
No total-language completeness result is obtained merely by writing down
such program pairs.

\paragraph{Verification, remaining gap, and outcome.}
The companion exact checks compare the factor-threshold predicate with the
prime-factor certificate criterion over small integers and every relevant
threshold, including primes, squares, repeated factors, and thresholds
outside $[2,N-1]$. They then implement the binary-search reduction and
recover the complete factorizations in those cases. Small examples also
test the prefix obstruction. Trial division is used only inside this finite
verification script; the polynomial verifier in Lemma 34.1 invokes the
established primality theorem, not trial division on binary integers.

\textbf{UNRESOLVED.} The concrete progress is a total language with fully
specified certificates on both sides and a proved deterministic equivalence
to factoring. The general search route has exact missing steps: totality
is not a polynomial search procedure, the complement of a prefix predicate
is not supplied by the opposite-answer verifier, and arbitrary verifier
pairs need not define total languages. A solution still requires either
one unconditional lower bound for a genuine two-sided-certificate language,
or a uniform decision method for every such language. No such bound or
method is established, and no novelty is claimed for the standard reductions
or elementary counterexamples.

\subsection{Sources}
\begin{itemize}
\item[S1]Scott Aaronson, PHYS771 Lecture 6: P, NP, and Friends.
\url{https://www.scottaaronson.com/democritus/lec6.html}.
\textit{Formulation source.}
\item[S2]Very few problems are in NP intersect coNP but not known to be in P. What to make of that?.
\url{https://blog.computationalcomplexity.org/2024/09/very-few-problems-are-in-np-intersect.html?m=0}.
\textit{Further reference; not a proof endorsement.}
\item[S3]Random Permutations in Computational Complexity.
\url{https://arxiv.org/abs/2511.08786}.
\textit{Further reference; not a proof endorsement.}
\item[E1]Manindra Agrawal, Neeraj Kayal, and Nitin Saxena,
\emph{PRIMES is in P}, Annals of Mathematics 160(2) (2004), pp. 781--793.
\url{https://annals.math.princeton.edu/2004/160-2/p12}.
\url{https://annals.math.princeton.edu/wp-content/uploads/annals-v160-n2-p12.pdf}.
\textit{Primary source for deterministic polynomial-time primality testing;
consulted 27 September 2026. Used to verify supplied prime factors, not
as an algorithm to discover a factorization.}
\item[E2]Nimrod Megiddo and Christos H. Papadimitriou,
\emph{On Total Functions, Existence Theorems and Computational Complexity},
Theoretical Computer Science 81(2) (1991), pp. 317--324, Section 1.
\url{https://theory.stanford.edu/~megiddo/pdf/papadimX.pdf}.
\textit{Primary definition of total polynomially verifiable search relations;
consulted 27 September 2026. No unproved search-class separation is used.}
\end{itemize}
\end{document}
